%%% This file contains definitions of various useful macros and environments %%%
%%% Please add more macros here instead of cluttering other files with them. %%%

%%% Switches based on thesis type

\def\TypeBc{bc}
\def\TypeMgr{mgr}
\def\TypePhD{phd}
\def\TypeRig{rig}

\ifx\ThesisType\TypeBc
\def\ThesisTypeName{bachelor}
\def\ThesisTypeTitle{BACHELOR THESIS}
\def\ThesisTypeCSGenitive{bakalářské}
\fi

\ifx\ThesisType\TypeMgr
\def\ThesisTypeName{master}
\def\ThesisTypeTitle{MASTER THESIS}
\def\ThesisTypeCSGenitive{diplomové}
\fi

\ifx\ThesisType\TypePhD
\def\ThesisTypeName{doctoral}
\def\ThesisTypeTitle{DOCTORAL THESIS}
\def\ThesisTypeCSGenitive{disertační}
\fi

\ifx\ThesisType\TypeRig
\def\ThesisTypeName{rigorosum}
\def\ThesisTypeTitle{RIGOROSUM THESIS}
\def\ThesisTypeCSGenitive{rigorozní}
\fi

\ifx\ThesisTypeName\undefined
\PackageError{thesis}{Unknown thesis type.}{Please check the definition of ThesisType in metadata.tex.}
\fi

%%% Switches based on study program language

\def\LangCS{cs}
\def\LangEN{en}

\ifx\StudyLanguage\LangCS
\else\ifx\StudyLanguage\LangEN
\else\PackageError{thesis}{Unknown study language.}{Please check the definition of StudyLanguage in metadata.tex.}
\fi\fi

%%% Minor tweaks of style

% These macros employ a little dirty trick to convince LaTeX to typeset
% chapter headings sanely, without lots of empty space above them.
% Feel free to ignore.
\makeatletter
\def\@makechapterhead#1{
  {\parindent \z@ \raggedright \normalfont
   \Huge\bfseries \thechapter\quad #1
   \par\nobreak
   \vskip 20\p@
}}
\def\@makeschapterhead#1{
  {\parindent \z@ \raggedright \normalfont
   \Huge\bfseries #1
   \par\nobreak
   \vskip 20\p@
}}
\makeatother

% This macro defines a chapter, which is not numbered, but is included
% in the table of contents.
\def\chapwithtoc#1{
\chapter*{#1}
\addcontentsline{toc}{chapter}{#1}
}

% Draw black "slugs" whenever a line overflows, so that we can spot it easily.
\overfullrule=1mm

%%% Macros for definitions, theorems, claims, examples, ... (requires amsthm package)

\theoremstyle{plain}
\newtheorem{thm}{Theorem}
\newtheorem{lemma}[thm]{Lemma}
\newtheorem{claim}[thm]{Claim}
\newtheorem{defn}{Definition}

\theoremstyle{remark}
\newtheorem*{cor}{Corollary}
\newtheorem*{rem}{Remark}
\newtheorem*{example}{Example}

%%% Style of captions of floating objects (figures etc.)

\ifcsname DeclareCaptionStyle\endcsname
\DeclareCaptionStyle{thesis}{style=base,font=small,labelfont=bf,labelsep=quad}
\captionsetup{style=thesis}
\captionsetup[algorithm]{style=thesis,singlelinecheck=off}
\captionsetup[listing]{style=thesis,singlelinecheck=off}
\fi

%%% An environment for typesetting of program code and input/output
%%% of programs.
\DefineVerbatimEnvironment{code}{Verbatim}{fontsize=\small, frame=single}

% Floating listings, used in the same way as the figure environment
\ifcsname DeclareNewFloatType\endcsname
\DeclareNewFloatType{listing}{}
\floatsetup[listing]{style=ruled}
\floatname{listing}{Program}
\fi

%%% The field of all real and natural numbers
\newcommand{\R}{\mathbb{R}}
\newcommand{\N}{\mathbb{N}}

%%% Useful operators for statistics and probability
\DeclareMathOperator{\pr}{\textsf{P}}
\DeclareMathOperator{\E}{\textsf{E}}
\DeclareMathOperator{\var}{\textrm{var}}
\DeclareMathOperator{\sd}{\textrm{sd}}

%%% Transposition of a vector/matrix
\newcommand{\T}[1]{#1^\top}

%%% Asymptotic "O"
\def\O{\mathcal{O}}

%%% Various math goodies
\newcommand{\goto}{\rightarrow}
\newcommand{\gotop}{\stackrel{P}{\longrightarrow}}
\newcommand{\maon}[1]{o(n^{#1})}
\newcommand{\abs}[1]{\left|{#1}\right|}
\newcommand{\dint}{\int_0^\tau\!\!\int_0^\tau}
\newcommand{\isqr}[1]{\frac{1}{\sqrt{#1}}}

%%% TODO items: remove before submitting :)
\newcommand{\xxx}[1]{\textcolor{red!}{#1}}

%%% Detailed settings of bibliography

\ifx\citet\undefined\else

% Maximum number of authors of a single work. If exceeded, "et al." is used.
%\ExecuteBibliographyOptions{maxnames=2}
% The same setting specific to citations using \citet{...}
\ExecuteBibliographyOptions{maxcitenames=2}
% The same settings specific to the list of literature
%\ExecuteBibliographyOptions{maxbibnames=2}

% Shortening first names of authors: "E. A. Poe" instead of "Edgar Allan Poe"
%\ExecuteBibliographyOptions{giveninits}
% The same without dots ("EA Poe")
%\ExecuteBibliographyOptions{terseinits}

% If your bibliography entries are hard to break into lines, try this mode:
%\ExecuteBibliographyOptions{block=ragged}

% Possibly reverse the names of the authors with the non-ISO styles:
%\DeclareNameAlias{default}{family-given}

% Use caps-and-small-caps for family names in ISO 690 style.
\let\familynameformat=\textsc

% We want to separate multiple authors in citations by commas
% (while we use semicolons in the bibliography as per the ISO standard)
\DeclareDelimFormat[textcite]{multinamedelim}{\addcomma\space}
\DeclareDelimFormat[textcite]{finalnamedelim}{\space and~}

\fi

%%%%%%%%%%%%%%%%%%%%%%%%%%%%%%%%%%%%%%%%%%%%%%%%%%%%%%
%%%%%%%%%%% Default listing format %%%%%%%%%%%%%%%%%%%

\definecolor{gscomment}{RGB}{13, 97, 4}
\definecolor{gsline}{RGB}{238,232,213}
\definecolor{gslineno}{RGB}{131,148,150}


\lstset{
    frame=l,
    linewidth=160mm,
    tabsize=4,
    numbers=left,
    numbersep=7pt,
    firstnumber=auto,
    numberstyle=\tiny\ttfamily\color{gslineno},
    rulecolor=\color{gsline},
    basicstyle=\footnotesize\ttfamily,
    commentstyle=\color{gscomment},
}

%%%%%%%%%%% Default listing format %%%%%%%%%%%%%%%%%%%
%%%%%%%%%%%%%%%%%%%%%%%%%%%%%%%%%%%%%%%%%%%%%%%%%%%%%%

%%%%%%%%%%%%%%%%%%%%%%%%%%%%%%%%%%%%%%%%%%%%%%%%%%%%%%
%%%%%%%%%%% YAML syntax highlighting %%%%%%%%%%%%%%%%%

% http://tex.stackexchange.com/questions/152829/how-can-i-highlight-yaml-code-in-a-pretty-way-with-listings

% here is a macro expanding to the name of the language
% (handy if you decide to change it further down the road)

\definecolor{sbase}{RGB}{17, 37, 64}

\newcommand\YAMLbasestyle{\ttfamily\color{sbase}\bfseries\footnotesize}
\newcommand\YAMLvaluestyle{\ttfamily\color{brown}\mdseries\footnotesize}

\makeatletter

\newcommand\language@yaml{yaml}

\expandafter\expandafter\expandafter\lstdefinelanguage
\expandafter{\language@yaml}
{
  keywords={true,false,null,y,n},
  keywordstyle=\color{darkgray}\bfseries,
  basicstyle=\YAMLbasestyle,                                 % assuming a key comes first
  sensitive=false,
  comment=[l]{\#},
  morecomment=[s]{/*}{*/},
  stringstyle=\YAMLvaluestyle\ttfamily,
  morestring=[b]',
  morestring=[b]",
  literate = {:}{{\YAMLbasestyle}{:}{\YAMLvaluestyle}}1
             {-\ }{{\YAMLbasestyle}{-\ }{\YAMLvaluestyle}}2
}

% switch to key style at EOL
\lst@AddToHook{EveryLine}{\ifx\lst@language\language@yaml\YAMLbasestyle\fi}
\makeatother

%%%%%%%%%%% YAML syntax highlighting %%%%%%%%%%%%%%%%%
%%%%%%%%%%%%%%%%%%%%%%%%%%%%%%%%%%%%%%%%%%%%%%%%%%%%%%

%%%%%%%%%%%%%%%%%%%%%%%%%%%%%%%%%%%%%%%%%%%%%%%%%%%%%%
%%%%%%%%%%% C# and C++ syntax highlighting %%%%%%%%%%%

% Edited version of https://gist.github.com/Cheesebaron/d840258a9c41619dbdf0

\definecolor{ssymbol}{RGB}{44,55,58}
\definecolor{skeyword}{RGB}{24, 78, 143}
\definecolor{sident}{RGB}{23, 143, 131}
\definecolor{sstring}{RGB}{207, 106, 19}
\definecolor{sdigit}{RGB}{211,54,130}
\definecolor{sfunction}{RGB}{153, 131, 20}
\definecolor{sclass}{RGB}{25, 138, 13}


\newcommand\digitstyle{\color{sdigit}}
\newcommand\symbolstyle{\color{ssymbol}}
\makeatletter

\newcommand\language@csharp{csharp}


\newcommand{\ProcessDigit}[1]
{%
  \ifnum\lst@mode=\lst@Pmode\relax%
   {\digitstyle #1}%
  \else
    #1%
  \fi
}

\expandafter\expandafter\expandafter\lstdefinelanguage
\expandafter{\language@csharp}
{
  language=[Sharp]C,
  morecomment=[s]{/*+}{*/},
  morecomment=[s]{/*-}{*/},
  morestring=[b]',
  morekeywords={  abstract, event, new, struct,
                as, explicit, null, switch,
                base, extern, object, this,
                bool, false, operator, throw,
                break, finally, out, true,
                byte, fixed, override, try,
                case, float, params, typeof,
                catch, for, private, uint,
                char, foreach, protected, ulong,
                checked, goto, public, unchecked,
                class, if, readonly, unsafe,
                const, implicit, ref, ushort,
                continue, in, return, using,
                decimal, int, sbyte, virtual,
                default, interface, sealed, volatile,
                delegate, internal, short, void, record,
                do, is, sizeof, while, required,
                double, lock, stackalloc,
                else, long, static, async, await,
                enum, namespace, string, var,
                template, typename, constexpr, % C++ keywords
                int8_t, int16_t, int32_t, int64_t, size_t,
                uint8_t, uint16_t, uint32_t, uint64_t,
                auto, unsigned, friend, inline},
  keywordstyle=\color{skeyword},
  showstringspaces=false,
  stringstyle=\color{sstring},
  identifierstyle=\color{sident},
  extendedchars=true,
  literate=
    % {0}{{{\ProcessDigit{0}}}}1
    % {1}{{{\ProcessDigit{1}}}}1
    % {2}{{{\ProcessDigit{2}}}}1
    % {3}{{{\ProcessDigit{3}}}}1
    % {4}{{{\ProcessDigit{4}}}}1
    % {5}{{{\ProcessDigit{5}}}}1
    % {6}{{{\ProcessDigit{6}}}}1
    % {7}{{{\ProcessDigit{7}}}}1
    % {8}{{{\ProcessDigit{8}}}}1
    % {9}{{{\ProcessDigit{9}}}}1
    {\}}{{\symbolstyle{\}}}}1
    {\{}{{\symbolstyle{\{}}}1
    {(}{{\symbolstyle{(}}}1
    {)}{{\symbolstyle{)}}}1,
    % {=}{{\symbolstyle{$=$}}}1
    % {;}{{\symbolstyle{$;$}}}1
    % {>}{{\symbolstyle{$>$}}}1
    % {<}{{\symbolstyle{$<$}}}1
    % {\%}{{\symbolstyle{$\%$}}}1,
  moredelim=[is][\color{sfunction}\bfseries]{\#}{\#},
  moredelim=[is][\color{sclass}\bfseries]{$}{$}
}
\makeatother

%%%%%%%%%%% C# syntax highlighting %%%%%%%%%%%%%%%%%%%
%%%%%%%%%%%%%%%%%%%%%%%%%%%%%%%%%%%%%%%%%%%%%%%%%%%%%%

%%%%%%%%%%%%%%%%%%%%%%%%%%%%%%%%%%%%%%%%%%%%%%%%%%%%%%
%%%%%%%%%%% Shell syntax highlighting %%%%%%%%%%%%%%%%

\definecolor{sexe}{RGB}{17, 37, 64}
\definecolor{sopt}{RGB}{153, 131, 20}
\definecolor{sarg}{RGB}{8, 143, 161}

\makeatletter

\newcommand\language@shell{shell}

\expandafter\expandafter\expandafter\lstdefinelanguage
\expandafter{\language@shell}
{
  basicstyle=\footnotesize\ttfamily\bfseries\color{sexe},
  moredelim=[s][\color{sopt}]{\ -}{\ },
  moredelim=[s][\color{sopt}]{\ --}{\ }
}

\makeatother


%%%%%%%%%%% Shell syntax highlighting %%%%%%%%%%%%%%%%
%%%%%%%%%%%%%%%%%%%%%%%%%%%%%%%%%%%%%%%%%%%%%%%%%%%%%%